\color{unimain}{\subsection{Логика отключения}}
\color{black}

\begin{enumerate}

\item
Функциональный блок логики отключения (ЛО) предназначен для подачи сигнала на отключение
силового выключателя. 

\item
Функциональная схема ЛО представлена на рисунке \ref{pict_lo:LO}.

\begin{figure}[H]
\centering
\includegraphics[width=0.8\textwidth,height=0.8\textheight,keepaspectratio]{\fbpath/08.01 СВ/761.9901 ЛО/img/LO.pdf}
\captionof{figure}{Функциональная схем ЛО}\label{pict_lo:LO}
\end{figure}

\item
Команда на отключение выключателя \textit{«Общая ЛО/ЛО: Отключить»} формируется:
\begin{itemize}
	\item при срабатывании внутренних защит устройства (сигнал \textit{«Срабатывание»});
	\item при срабатывании ЗНФ (сигнал \textit{«ЗНФ СВ/ЗНФ: Срабатывание»});
	\item при действии УРОВ на отключение «своего» выключателя (сигнал \textit{«УРОВ СВ/УРОВ: Срабатывание на себя»});
	\item от сигнала об аварийном давлении элегаза в ТТ, если введен программный ключ \textit{«КСВ СВ/Контроль ТТ: Откл. от Авар. давл. SF6 ТТ»};
	\item от команды оперативного отключения выключателя (сигнал \textit{«БКУ СВ/БКУ: РКО»});
	\item от сигналов внешнего отключения выключателя (сигнал \textit{«ЛО: Внеш. откл от РЗА 1СШ»} и \textit{«ЛО: Внеш. откл от РЗА 2СШ»}).
\end{itemize}

\item
Ввод в работу логики отключения выключателя осуществляется с помощью программного ключа «ЛО».
При необходимости может быть введен подхват команды на отключение от токовых реле, установленных в
цепях ЭМО, в этом случае сигнал на отключение удерживается в сработанном состоянии до исчезновения тока
через ЭМО (пока блок-контакты выключателя не разорвут цепи отключения). Ввод осуществляется
программным ключом \textbf{«Подхват откл от РТ ЭМО»}. 

\item
При наличии сигнала \textbf{«КСВ СВ/ Блокировка упр.: Блок. отключения»}, формируемого функциональным блоком контроля силового
выключателя, команды на отключение выключателя блокируется.

\item
Выдержка времени на возврат \textbf{«Тпрод откл»} задает время продления сигнала отключения. 

\item
Оперативный ввод/вывод ЛО производится с помощью БУ (блоков управления функциями) \textit{«ЛО/ЛО:Вывод»}.

\item
Параметры для настройки функции ЛО представлены в таблице \ref{app_set:lo}.

\end{enumerate}